\chapter*{Preface}

\addcontentsline{toc}{chapter}{Preface}

The retina is the only part of the central nervous system that can be directly visualized non-invasively. Owing to its shared embryological origin, anatomical structure, and physiological function with the brain, retinal microvascular alterations have been associated with various systemic diseases. Advances in retinal imaging and artificial intelligence have transformed the retina from an organ primarily examined for ocular diseases into a valuable source of quantitative biomarkers for systemic health. Recent advances in artificial intelligence (AI) have further enhanced the potential of retinal imaging as a non-invasive biomarker for cerebrovascular health. These developments have transformed retinal image analysis from a purely qualitative clinical assessment to large-scale, quantitative, and automated evaluation of retinal images, enabling early detection of disease, risk stratification, and monitoring of disease progression.

The motivation for this handbook originated from my own experience as a Computer Science PhD student entering the field of retinal imaging. Although building machine learning models was familiar, understanding biomedical terminology, retinal anatomy, clinical concepts, and disease mechanisms proved to be a much greater challenge than expected. Many important concepts were scattered across ophthalmology textbooks, neuroscience references, epidemiology papers, and artificial intelligence literature. This fragmentation made it difficult to build a coherent understanding of the field. This handbook was therefore developed as a living knowledge base that helps translate complex clinical concepts into computational understanding. It is intended to evolve alongside scientific advances, serving as a ling knowledge base for researchers, clinicians, and students exploring the intersection of retinal imaging, vascular biology, and artificial intelligence. The handbook is organized into chapters, each focusing on a specific topic, with a glossary of terms provided for reference. The content is written in English, with Vietnamese translations provided for selected glossary terms to facilitate understanding for Vietnamese-speaking readers. Rather than providing definitive answers, this handbook aims to support continuous learning, critical thinking, and interdisciplinary research throughout the rapidly evolving field of retinal imaging and artificial intelligence.